\documentclass[10pt,a4paper]{amsart}
\usepackage[margin=19mm]{geometry}
\usepackage{amsmath,amssymb,mathtools}
\usepackage[scr=rsfs,cal=euler]{mathalfa}
\usepackage{xfrac}
\usepackage[unicode,hidelinks,bookmarks=false]{hyperref}
\newcommand{\ie}{i.e.\ }
\newcommand{\cf}{cf.\ }
\newcommand{\resp}{respectively\ }
\newcommand{\Subsection}{\subsection}
\providecommand{\nobreakdash}{\nobreak\hbox{-}}
\setlength{\emergencystretch}{2em}
\multlinegap0pt
\usepackage{bm}
\usepackage{bbm}
\usepackage{dsfont}
\usepackage{xspace}
\usepackage{cancel}
\usepackage{mathtools}
\usepackage{amsthm}
\makeatletter
\@ifundefined{theo}{\newtheorem{theo}{Theo}}{}
\@ifundefined{prop}{\newtheorem{prop}[theo]{Proposition}}{}
\makeatother
\title[Rectangular diagrams of foliations]{Rectangular diagrams of foliations}
\author{Mikhail Chernavskikh, Ivan Dynnikov}
\date{Source: arXiv:2508.07732v1 (CC BY 4.0)}
\begin{document}
\maketitle

\noindent\emph{Source and licence.} Rectangular diagrams of foliations. Version arXiv:2508.07732v1 (CC BY 4.0), distributed under the Creative Commons Attribution 4.0 International licence, as recorded in the arXiv licence metadata for this exact version and shown on the abstract page https://arxiv.org/abs/2508.07732v1. Full rights notice: licenses/arxiv-2508.07732-CC-BY-4.0.txt.\par\smallskip

\noindent\textit{Editorial scope.} Counted unit: the Prop whose source label is \texttt{normal=>rect-prop} in the article (source0.tex lines 1843--1848, source label \texttt{normal=>rect-prop}), delivered together with the article's own complete original proof of it (source0.tex lines 1850--1910, one \texttt{proof} environment of 61 lines). No mathematics is written, repaired or re-derived: statement and proof are the article's own text line for line. Required context retained uncounted: none. Every definition, display and label of those lines is retained unchanged. Omitted from this package and not redistributed: 1--1842, 1849--1849, 1911--2835. Nothing inside a retained excerpt is omitted or altered: every retained body is delivered exactly as the article's lines are written, so no omission receipt of any kind accompanies this package. Citations: the retained excerpts carry no citation command at all, so no bibliography entry of the article is transcribed and no reference is redistributed. Reference adaptation: none was needed and none was made; every reference inside the delivered unit resolves inside it, so no unresolved reference or citation is produced. Printed slips kept literal: no printed word, symbol, hypothesis or number of the counted statement or of its proof was altered. Layout and class: the article's own class is \texttt{amsart} with packages including amsmath, amssymb, mathrsfs, bbm, MnSymbol, longtable, mathtools, amsthm, rotating, epsfig, color, graphicx, hyperref; the wrapper is the lane's pinned \texttt{amsart} editorial wrapper at 10pt on A4, which restates the theorem-like environments the retained text needs and adds the article's own macro definitions that the retained text needs (none), with extra wrapper packages (bm, bbm, dsfont, xspace, cancel, mathtools). The delivered numbering is the unit's own. Assets: no retained span contains a figure environment, a graphics-inclusion command, a PDF-page inclusion, a table or a TikZ picture; no image file is copied into this package, the package declares no asset dependency and no image file is redistributed with it. Licence: version arXiv:2508.07732v1 of this article is distributed under the Creative Commons Attribution 4.0 International licence, as recorded in the arXiv licence metadata for this exact version and shown on the abstract page https://arxiv.org/abs/2508.07732v1. A full rights notice is delivered at licenses/arxiv-2508.07732-CC-BY-4.0.txt. Characters: every retained excerpt is pure ASCII, so the delivered engine, XeLaTeX with its default Latin Modern text font, reports no missing character.
\begin{prop}\label{normal=>rect-prop}
Let~$X,Y\subset\mathbb S^1$ be finite subsets with~$|X|,|Y|\geqslant3$, and let~$F\subset\mathbb S^3$
be a normal surface with respect to~$T(X,Y)$ such that every connected component of~$F$
has a non-empty intersection with both circles~$\mathbb S^1_{\tau=0}$ and~$\mathbb S^1_{\tau=1}$. Then~$F$
is normal isotopic to a surface of the form~$\widehat\Pi$, where~$\Pi$ is a rectangular diagram of a surface.
\end{prop}

\begin{proof}
We call a normal disc~$d$ of~$F$ \emph{vertical} if~$d$ intersects both circles~$\mathbb S^1_{\tau=0}$
and~$\mathbb S^1_{\tau=1}$. One can see that this can only happen when~$d$ is a quadrilateral,
and the two corners of~$d$ that lie on these circles are opposite to each other.

Let~$d_1,d_2,d_3,\ldots$ be all vertical normal discs of~$F$ (there may be an infinite number of them), and, for each~$i$,
let~$\alpha_i$ be a proper arc in~$d_i$ having endpoints at~$d_i\cap\mathbb S^1_{\tau=0}$ and~$d_i\cap\mathbb S^1_{\tau=1}$. We may assume without loss of generality that the surface~$F$
is transverse to all arcs of the form~$\widehat v$, where~$v\in\mathbb T^2$, outside the union~$\bigcup_i\alpha_i$,
whereas the arcs~$\alpha_i$ have the form~$\widehat v$.

Indeed, let us endow each $3$-simplex of~$T(X,Y)$ with an affine structure so that each arc of the form~$\widehat v$,
where~$v\in\mathbb T^2$, is a straight line segment. Then~$F$ is normal isotopic to a surface~$F'$ every
normal disc in which is either an affine triangle or a quadrilateral composed of two affine triangles in a $3$-simplex of~$T(X,Y)$,
where the normal isotopy can be chosen to be fixed on the $1$-skeleton of~$T(X,Y)$.
Moreover, vertical normal discs of~$F'$ can be chosen
in the form of the union of two affine triangles sharing a side of the form~$\widehat v$ for some~$v\in\mathbb T^2$.
The surface~$F'$ is transverse to the arcs of the form~$\widehat v$ everywhere except at a union of arcs
that can be taken for~$\alpha_i$.

In what follows we assume that~$\alpha_i=\widehat v_i$ for some~$v_i\in\mathbb T^2$, and~$F\setminus\bigcup_i\alpha_i$
is transverse to all arcs of the form~$\widehat v$.

Let~$V$ be a connected component of~$F\setminus\bigcup_i\alpha_i$. Define a map~$\omega:V\rightarrow\mathbb T^2$
by demanding that~$\widehat{\omega(p)}\cap V=p$ for all~$p\in V$. The pullpack of the standard locally
Euclidean metric on~$\mathbb T^2=\mathbb S^1\times\mathbb S^1$
is a locally Euclidean metric on~$V$. Let~$\widetilde V$ be the minimal compactification of~$V$
making this metric complete and keeping it locally Euclidean, and extend the map~$\omega$ to~$\widetilde V$ by continuity.
Clearly, $\omega$ is a locally isometric immersion~$\widetilde V\rightarrow\mathbb T^2$.

For every $3$-simplex~$\Delta$ of~$T(X,Y)$,
the closure of each connected component of the intersection~$V\cap\Delta$ is either a non-vertical normal disc
or one half of a vertical quadrilateral. This means that the image of~$V\cap\Delta$ under~$\omega$
has the form~$I\times J$, where each of~$I$ and~$J$ is either a closed or half-open interval of~$\mathbb S^1$.
Thus, the intersection of~$V$ with the $2$-skeleton of~$T(X,Y)$ cuts~$V$ into pieces isometric to such
subsets of~$\mathbb T^2$.

This implies that~$\widetilde V$ can be cut into finitely many rectangles, and the angle at any breaking point
of~$\partial\widetilde V$ is~$\pi/2$. Therefore, $\widetilde V$
is isometric either to a rectangle in the Euclidean plane, or a flat annulus with geodesic
boundary, or a flat two-torus. Now we claim that~$\omega$ is actually an embedding, and the latter
two cases (an annulus and a two-torus) are impossible.

Indeed, the vertical and horizontal arcs in~$\omega(\partial\widetilde V)$ come from the intersections of~$\overline V$
with~$\mathbb S^1_{\tau=0}$ and~$\mathbb S^1_{\tau=1}$, respectively. The equality~$\partial\widetilde V=\varnothing$
would mean that~$\overline V\cap(\mathbb S^1_{\tau=0}\cup\mathbb S^1_{\tau=1})=\varnothing$,
which implies~$\partial\overline V=\varnothing$. So, $V=\overline V$ would be a connected component of~$F$
disjoint from~$\mathbb S^1_{\tau=0}\cup\mathbb S^1_{\tau=1}$, which is assumed not to be the case.

If the entire~$\omega(\partial\widetilde V)$ is horizontal (respectively, vertical) in~$\mathbb T^2$,
then~$\overline V\cap\mathbb S^1_{\tau=1}=\varnothing$ (respectively, $\overline V\cap\mathbb S^1_{\tau=0}=\varnothing$),
which implies that~$V=\overline V$ is a spherical component of~$F$ intersecting only one of the
circles~$\mathbb S^1_{\tau=0}$ and~$\mathbb S^1_{\tau=1}$, and this again contradicts the hypothesis of the proposition.

Thus, $\widetilde V$ is a rectangle. The length of the sides of~$\widetilde V$ corresponding to
the intersections of~$\overline V$ with~$\mathbb S^1_{\tau=0}$ (respectively, $\mathbb S^1_{\tau=1}$) is equal to the variation of the
$\theta$-coordinate (respectively, $\varphi$-coordinate) on the boundary of a small neighborhood of such an intersection in~$V$, which implies that it is smaller than~$2\pi$. Therefore, $\omega$ is an embedding, and~$\omega(\widetilde V)$ is a rectangle in~$\mathbb T^2$.
Denote it by~$r(V)$.

Now let~$\Pi$ be the set of all rectangles of the form~$r(V)$, where~$V$ runs over the set of connected
components of~$F\setminus\bigcup_i\alpha_i$. One can see that the surface~$\widehat\Pi$ is normal parallel to~$F$.
\end{proof}

\end{document}
